% Elliptic Functions and Liouville's Theorems  (source §4, major expansion:
% the field of elliptic functions, no nonconstant holomorphic ones, and the
% three Liouville theorems -- residue sum, the valence/counting theorem, and the
% Abel relation sum(zeros) - sum(poles) in Lambda -- by contour integration over
% a period parallelogram)
\chapter{Elliptic Functions and Liouville's Theorems}
\label{ch:elliptic_liouville}

By \cref{rem:lat-translation-holo}, meromorphic functions on the torus
\(T = \C/\Lambda\) are the same as \(\Lambda\)-periodic meromorphic functions on
\(\C\) --- the \emph{elliptic functions}. This chapter establishes the rigid
constraints these functions obey. There are no nonconstant holomorphic ones, so
a nonconstant elliptic function must have poles; and the location and number of
its zeros and poles are governed by three theorems of Liouville, all proved by
integrating around a single period parallelogram. The third of them,
\(\sum(\text{zeros}) - \sum(\text{poles}) \in \Lambda\), is the analytic engine
of the entire group law: it returns in \cref{ch:abel} as Abel's sum law and is
what ultimately forces \(P + Q + R = O\) for collinear points in
\cref{thm:main}.

\section{The Field of Elliptic Functions}
\label{sec:ell-field}

\begin{definition}
\label{def:ell-elliptic}
A meromorphic function \(f\) on \(\C\) is \emph{elliptic} with respect to
\(\Lambda\), or \(\Lambda\)-periodic, if
\[
  f(z + \lambda) = f(z) \qquad \text{for all } \lambda \in \Lambda .
\]
Equivalently, \(f(z+\omega_1) = f(z)\) and \(f(z+\omega_2) = f(z)\) for a basis
\(\omega_1, \omega_2\).
\end{definition}

An elliptic function depends only on the class \([z]\), so it descends to a
meromorphic function \(T \to \CP^1\); conversely every meromorphic function on
\(T\) pulls back along \(\pi\) to a \(\Lambda\)-periodic one. The two viewpoints
are interchangeable, and we pass between them freely:
\[
  \{\text{meromorphic functions on } T\}
  \;\longleftrightarrow\;
  \{\Lambda\text{-elliptic functions on } \C\}.
\]

\begin{proposition}
\label{prop:ell-field}
The \(\Lambda\)-elliptic functions form a field, closed under differentiation:
if \(f, g\) are elliptic then so are \(f \pm g\), \(fg\), \(f/g\) (for
\(g \not\equiv 0\)), and \(f'\).
\end{proposition}

\begin{proof}
Sums, products, and quotients of meromorphic functions are meromorphic, and each
inherits \(\Lambda\)-periodicity pointwise, e.g.\ \((fg)(z+\lambda) =
f(z+\lambda)g(z+\lambda) = f(z)g(z)\). Differentiating \(f(z+\lambda)=f(z)\) in
\(z\) gives \(f'(z+\lambda) = f'(z)\), so \(f'\) is elliptic.
\end{proof}

The central problem of the theory is to describe this field. The answer,
reached in \cref{ch:weierstrass_wp}, is that it is generated by the Weierstrass
function \(\wp\) and its derivative \(\wp'\). Everything in this chapter is a
constraint that the generators must satisfy.

\section{No Nonconstant Holomorphic Elliptic Functions}
\label{sec:ell-liouville-const}

\begin{theorem}[Liouville's first theorem]
\label{thm:ell-const}
An elliptic function with no poles is constant.
\end{theorem}

\begin{proof}
Let \(f\) be elliptic and holomorphic on all of \(\C\). The closed fundamental
parallelogram
\(\overline{P} = \{ s\omega_1 + t\omega_2 : 0 \le s, t \le 1 \}\) is compact, so
the continuous function \(|f|\) attains a maximum \(M\) on \(\overline{P}\). By
periodicity every \(z \in \C\) is congruent to a point of \(\overline P\) and
\(f\) takes the same value there, so \(|f(z)| \le M\) for all \(z\): the entire
function \(f\) is bounded. By Liouville's theorem (a bounded entire function is
constant) \(f\) is constant. Equivalently, \(T\) is compact, so a holomorphic
\(f : T \to \C\) attains its maximum modulus and is constant by the maximum
modulus principle.
\end{proof}

Thus a nonconstant elliptic function must have at least one pole. The
quantitative versions of this statement are Liouville's remaining theorems.

\section{Integration Around a Period Parallelogram}
\label{sec:ell-contour}

All three theorems rest on one device. Fix a basis \(\omega_1, \omega_2\) and a
base point \(a \in \C\), and let
\[
  P_a = \{\, a + s\omega_1 + t\omega_2 : 0 \le s, t \le 1 \,\}
\]
be the translated fundamental parallelogram, with boundary \(\partial P_a\)
traversed counterclockwise. Because a nonconstant elliptic function has only
finitely many zeros and poles modulo \(\Lambda\), we may always choose \(a\) so
that \(\partial P_a\) passes through none of them
(\cref{fig:periodcontour}). Label the four sides
\[
  \mathrm{I}: a \to a{+}\omega_1, \quad
  \mathrm{II}: a{+}\omega_1 \to a{+}\omega_1{+}\omega_2, \quad
  \mathrm{III}: a{+}\omega_1{+}\omega_2 \to a{+}\omega_2, \quad
  \mathrm{IV}: a{+}\omega_2 \to a .
\]

\begin{figure}[ht]
\centering
\begin{tikzpicture}[scale=1.25, every node/.style={font=\small}]
  \coordinate (A)  at (0,0);
  \coordinate (B)  at (2.6,0);       % a + w1
  \coordinate (C)  at (3.5,2.2);     % a + w1 + w2
  \coordinate (D)  at (0.9,2.2);     % a + w2
  \fill[formteal!6] (A) -- (B) -- (C) -- (D) -- cycle;
  % side I (bottom) and III (top): direction-matched single arrows
  \draw[postaction={decorate, decoration={markings, mark=at position 0.55 with {\arrow[draw=formrust!75!black]{>}}}}, formrust!80!black, line width=0.9pt] (A) -- (B);
  \draw[postaction={decorate, decoration={markings, mark=at position 0.55 with {\arrow[draw=formrust!75!black]{>}}}}, formrust!80!black, line width=0.9pt] (D) -- (C);
  % side II (right) and IV (left): double arrows
  \draw[postaction={decorate, decoration={markings, mark=at position 0.5 with {\arrow[draw=formgold!75!black]{>}}, mark=at position 0.62 with {\arrow[draw=formgold!75!black]{>}}}}, formgold!80!black, line width=0.9pt] (B) -- (C);
  \draw[postaction={decorate, decoration={markings, mark=at position 0.5 with {\arrow[draw=formgold!75!black]{>}}, mark=at position 0.62 with {\arrow[draw=formgold!75!black]{>}}}}, formgold!80!black, line width=0.9pt] (A) -- (D);
  % vertices
  \foreach \p in {A,B,C,D}{\fill[black] (\p) circle (1.2pt);}
  \node[below left]  at (A) {$a$};
  \node[below right] at (B) {$a+\omega_1$};
  \node[above right] at (C) {$a+\omega_1+\omega_2$};
  \node[above left]  at (D) {$a+\omega_2$};
  \node[formrust!70!black] at (1.3,-0.22) {I};
  \node[formrust!70!black] at (2.2,2.42) {III};
  \node[formgold!60!black] at (3.25,1.1) {II};
  \node[formgold!60!black] at (0.32,1.1) {IV};
  % some zeros (open) and poles (cross) inside
  \draw[formteal!70!black, line width=0.7pt] (1.5,0.8) circle (2.4pt);
  \draw[formteal!70!black, line width=0.7pt] (2.3,1.4) circle (2.4pt);
  \node[formplum] at (1.05,1.45) {$\times$};
  \node[formplum] at (2.05,0.55) {$\times$};
  \node[font=\scriptsize, anchor=west, align=left] at (3.9,1.1)
    {$\circ$ zeros\\ $\times$ poles};
\end{tikzpicture}
\caption{A period parallelogram \(P_a\), chosen so its boundary avoids the zeros
(\(\circ\)) and poles (\(\times\)) of \(f\). For an elliptic integrand the
integral over side III cancels that over side I (they are translates by
\(\omega_2\) traversed oppositely), and II cancels IV; the boundary integral
vanishes.}
\label{fig:periodcontour}
\end{figure}

\begin{lemma}[Opposite sides cancel]
\label{lem:ell-cancel}
If \(g\) is elliptic and has no poles on \(\partial P_a\), then
\(\oint_{\partial P_a} g(z)\,dz = 0\).
\end{lemma}

\begin{proof}
Side \(\mathrm{III}\) is the translate of side \(\mathrm{I}\) by \(\omega_2\),
traversed in the opposite direction. Substituting \(z = w + \omega_2\) and using
\(g(w+\omega_2) = g(w)\),
\[
  \int_{\mathrm{III}} g\,dz
  = -\int_{\mathrm{I}} g(w+\omega_2)\,dw
  = -\int_{\mathrm{I}} g(w)\,dw ,
\]
so \(\int_{\mathrm{I}} + \int_{\mathrm{III}} = 0\). Likewise side
\(\mathrm{II}\) is the translate of side \(\mathrm{IV}\) by \(\omega_1\), traversed
oppositely, so \(\int_{\mathrm{II}} + \int_{\mathrm{IV}} = 0\). Adding gives
\(\oint_{\partial P_a} g\,dz = 0\).
\end{proof}

\section{Liouville's Theorems}
\label{sec:ell-liouville}

\subsection{The sum of residues vanishes}

\begin{theorem}[Liouville's second theorem: residues]
\label{thm:ell-residues}
For a nonconstant elliptic function \(f\), the sum of the residues over one
period parallelogram is zero:
\[
  \sum_{[p] \in T} \Res_{p} f = 0 .
\]
\end{theorem}

\begin{proof}
Choose \(a\) so that \(\partial P_a\) avoids the poles of \(f\). The interior of
\(P_a\) contains exactly one representative of each pole class. By the residue
theorem and \cref{lem:ell-cancel} applied to \(g = f\),
\[
  2\pi i \sum_{p \in P_a} \Res_p f
  = \oint_{\partial P_a} f(z)\,dz = 0 . \qedhere
\]
\end{proof}

\begin{corollary}
\label{cor:ell-order2}
No elliptic function has a single simple pole in a period parallelogram. The
\emph{order} of a nonconstant elliptic function --- the number of poles, counted
with multiplicity, in one period parallelogram --- is at least \(2\).
\end{corollary}

\begin{proof}
A function with one simple pole and no others would have a single nonzero
residue, contradicting \cref{thm:ell-residues}. So either there are at least two
poles, or a single pole of order \(\ge 2\); in both cases the order is at least
\(2\).
\end{proof}

This is the first hint of the Weierstrass function: the smallest a nonconstant
elliptic function can be is order two, realized by a single double pole. That
function is \(\wp\) (\cref{ch:weierstrass_wp}).

\subsection{The counting theorem}

\begin{theorem}[Liouville's third theorem: valence]
\label{thm:ell-valence}
A nonconstant elliptic function \(f\) of order \(N\) has exactly \(N\) zeros and
\(N\) poles in each period parallelogram, counted with multiplicity. More
generally, for every \(c \in \CP^1\) the equation \(f(z) = c\) has exactly \(N\)
solutions there; \(f : T \to \CP^1\) is an \(N\)-to-one branched cover.
\end{theorem}

\begin{proof}
The logarithmic derivative \(f'/f\) is elliptic by \cref{prop:ell-field}, with a
simple pole at each zero and pole of \(f\); its residue there is
\(\Res_p(f'/f) = \ord_p(f)\). Choosing \(a\) so \(\partial P_a\) avoids these
points, \cref{lem:ell-cancel} and the residue theorem give
\[
  0 = \oint_{\partial P_a} \frac{f'}{f}\,dz
  = 2\pi i \sum_{p \in P_a} \ord_p(f)
  = 2\pi i \bigl( Z - P \bigr),
\]
where \(Z\) and \(P\) are the numbers of zeros and poles with multiplicity.
Hence \(Z = P = N\). Applying this to \(f - c\) (also elliptic, of order \(N\),
with the same poles as \(f\)) shows \(f - c\) has \(N\) zeros, i.e.\ \(f\) takes
the value \(c\) exactly \(N\) times; the case \(c = \infty\) is the pole count.
\end{proof}

\subsection{The Abel relation}

The previous theorems used elliptic integrands. The decisive third statement
uses \(z\,f'(z)/f(z)\), whose factor \(z\) is \emph{not} periodic; the failure of
cancellation is precisely a lattice vector.

\begin{theorem}[Liouville's fourth theorem: the Abel relation]
\label{thm:ell-abel}
Let \(f\) be a nonconstant elliptic function. Listing its zeros \(z_1, \dots,
z_N\) and poles \(p_1, \dots, p_N\) in a period parallelogram with multiplicity,
\[
  \sum_{k=1}^{N} z_k \;-\; \sum_{k=1}^{N} p_k \;\in\; \Lambda .
\]
Equivalently, in the group \(T\), \(\sum_k [z_k] = \sum_k [p_k]\).
\end{theorem}

\begin{proof}
Consider \(g(z) = z\,f'(z)/f(z)\). At a zero or pole \(p\) of \(f\) of order
\(m = \ord_p(f)\) one has \(f'/f \sim m/(z-p)\), so \(\Res_p g = m\,p\). With
\(\partial P_a\) avoiding the zeros and poles, the residue theorem gives
\[
  \oint_{\partial P_a} g(z)\,dz
  = 2\pi i \sum_{p \in P_a} \ord_p(f)\, p
  = 2\pi i\Bigl( \sum_k z_k - \sum_k p_k \Bigr).
\tag{$\ast$}
\]
Now evaluate the contour integral directly, pairing opposite sides. On side
\(\mathrm{III}\) put \(z = w + \omega_2\) with \(w\) running along
\(\mathrm{I}\) reversed; since \(f'/f\) is \(\omega_2\)-periodic,
\[
  \int_{\mathrm{I}} + \int_{\mathrm{III}}
  = \int_{\mathrm{I}} \bigl[\, z - (z+\omega_2) \,\bigr]\frac{f'}{f}\,dz
  = -\,\omega_2 \int_{a}^{a+\omega_1} \frac{f'}{f}\,dz .
\]
The remaining integral is the change in \(\log f\) along side \(\mathrm{I}\);
because \(f(a+\omega_1) = f(a)\), this change is \(2\pi i\,k_1\) for some integer
\(k_1\) (the winding number of \(f\) along the side). Thus
\(\int_{\mathrm{I}} + \int_{\mathrm{III}} = -2\pi i\,k_1\,\omega_2\). Symmetrically,
pairing \(\mathrm{II}\) with \(\mathrm{IV}\) and using
\(\omega_1\)-periodicity,
\[
  \int_{\mathrm{II}} + \int_{\mathrm{IV}}
  = \omega_1 \int_{a}^{a+\omega_2} \frac{f'}{f}\,dz
  = 2\pi i\,k_2\,\omega_1
\]
for an integer \(k_2\). Adding,
\[
  \oint_{\partial P_a} g\,dz = 2\pi i\,( k_2\,\omega_1 - k_1\,\omega_2 ).
\tag{$\ast\ast$}
\]
Comparing \((\ast)\) and \((\ast\ast)\) and cancelling \(2\pi i\),
\[
  \sum_k z_k - \sum_k p_k = k_2\,\omega_1 - k_1\,\omega_2 \in \Lambda . \qedhere
\]
\end{proof}

\begin{keyidea}{The relation behind the group law}
\Cref{thm:ell-abel} says: for any elliptic function, the zeros and poles, added
as points of \(T\), are equal. Read backward it is a recipe --- if a divisor of
points has zero degree and its points sum to \(0\) in \(T\), it is the divisor of
an elliptic function (\cref{ch:riemann_roch} supplies the converse). When the
points are the intersections of the cubic with a line, this congruence becomes
\(P + Q + R = O\). The chord--tangent law is \cref{thm:ell-abel} in disguise.
\end{keyidea}

\section{Consequences and the Road Ahead}
\label{sec:ell-consequences}

The four theorems pin down the coarse shape of every elliptic function:

\begin{itemize}
\item it has no poles only if it is constant (\cref{thm:ell-const});
\item its residues sum to zero, so it cannot have order one
  (\cref{thm:ell-residues,cor:ell-order2});
\item it attains every value the same number \(N \ge 2\) of times
  (\cref{thm:ell-valence});
\item its zeros and poles satisfy \(\sum z_k \equiv \sum p_k \pmod \Lambda\)
  (\cref{thm:ell-abel}).
\end{itemize}

The minimal case \(N = 2\) is therefore the place to look for a generator of the
field. A function with a single double pole at \([0]\) would, by
\cref{thm:ell-valence}, have two zeros, and by \cref{thm:ell-abel} those zeros
\(z_1, z_2\) would satisfy \(z_1 + z_2 \equiv 0\), i.e.\ \(z_2 \equiv -z_1\):
the zeros are antipodal. This is exactly the behaviour of \(\wp - c\), and it is
why \(\wp\) is an even function. Whether such a function exists at all is a
dimension count, answered by Riemann--Roch in \cref{ch:riemann_roch}; its
explicit construction is \cref{ch:weierstrass_wp}.

\section{Summary}
\label{sec:ell-summary}

Elliptic functions are the meromorphic functions on \(T\); they form a
differential field. A nonconstant one must have poles (\cref{thm:ell-const}), of
total order \(N \ge 2\) (\cref{cor:ell-order2}), and it takes every value exactly
\(N\) times (\cref{thm:ell-valence}). Integrating \(f\), \(f'/f\), and
\(z\,f'/f\) around a period parallelogram yields, respectively, the vanishing of
the residue sum, the equality of zero and pole counts, and the Abel relation
\(\sum(\text{zeros}) - \sum(\text{poles}) \in \Lambda\)
(\cref{thm:ell-residues,thm:ell-valence,thm:ell-abel}). The last is the analytic
heart of the group law and reappears as Abel's sum law in \cref{ch:abel}.

\section*{Exercises}
\addcontentsline{toc}{section}{Exercises}

\begin{exercise}
\label{ex:ell-field}
Verify in detail that the \(\Lambda\)-elliptic functions form a field closed
under \(d/dz\) (\cref{prop:ell-field}). Show in addition that if \(f\) is
elliptic then so is \(f(z+a)\) for any fixed \(a \in \C\), and that an even
elliptic function (\(f(-z) = f(z)\)) has \(f'\) odd.
\end{exercise}

\begin{exercise}
\label{ex:ell-no-order-one}
Give a complete proof that there is no elliptic function of order one. Then show
that an elliptic function of order two has either two distinct simple poles with
opposite residues, or a single double pole with residue zero.
\end{exercise}

\begin{exercise}
\label{ex:ell-double-pole-zeros}
Suppose \(f\) is elliptic of order two with a single double pole at \([0]\).
Using \cref{thm:ell-valence,thm:ell-abel}, show that the two zeros of \(f\) are
of the form \([w]\) and \([-w]\) for some \(w\). Conclude that \(f(z) - f(z_0)\)
has zeros at \(\pm z_0\) for generic \(z_0\), the signature of an even function.
\end{exercise}

\begin{exercise}
\label{ex:ell-sum-values}
Let \(f\) be elliptic of order \(N\) and let \(c_1, c_2 \in \C\) be two values.
Using \cref{thm:ell-abel} applied to \(f - c_1\) and to \(f - c_2\), show that
the sum of the \(N\) solutions of \(f(z) = c_1\) equals, modulo \(\Lambda\), the
sum of the \(N\) solutions of \(f(z) = c_2\). (The common value is an invariant
of \(f\).)
\end{exercise}

\begin{exercise}
\label{ex:ell-winding}
In the proof of \cref{thm:ell-abel} the integers \(k_1, k_2\) are winding
numbers of \(f\) along the two sides of \(P_a\). Show directly that
\(\int_a^{a+\omega_1} (f'/f)\,dz \in 2\pi i\,\Z\) whenever \(f(a+\omega_1) =
f(a)\) and \(f\) has no zeros or poles on the segment, and interpret \(k_1\) as
the number of times \(f\) winds around \(0\).
\end{exercise}
